\documentclass[conference,comsoc,twoside,caption=false]{IEEEtran}
\usepackage{newtxtext,newtxmath}
\ifCLASSOPTIONcompsoc
  \usepackage[caption=false,font=footnotesize,labelfont=sf,textfont=sf]{subfig}
\else
  \usepackage[caption=false,font=footnotesize,labelfont=small]{subfig}
\fi
\captionsetup[figure]{font=footnotesize}
\ifCLASSOPTIONcaptionsoff
  \usepackage[nomarkers]{endfloat}
  \let\MYoriglatexcaption\caption
  \renewcommand{\caption}[2][\relax]{\MYoriglatexcaption[#2]{#2}}
\fi
\usepackage[utf8]{inputenc}
\usepackage{ragged2e}
\usepackage{graphicx}
\usepackage{tabulary}
\usepackage{booktabs}
\usepackage{mathtools}
\let\openbox\relax
\usepackage{amsmath}
\usepackage{amsthm}
\let\Bbbk\relax
\usepackage{amssymb}
\usepackage{amsfonts}
\interdisplaylinepenalty=2500
\usepackage[noadjust]{cite}
\usepackage{balance}
\usepackage{tabularx}
\usepackage{comment}
\usepackage{setspace}
\usepackage{xspace}
\usepackage{acronym}
\usepackage{multirow}
\usepackage{upgreek}
\usepackage{nicefrac}
\usepackage{ifthen}
\usepackage[bookmarks=false]{hyperref}
\hypersetup{draft}
\usepackage{caption}
\usepackage{cleveref}
\usepackage[inline]{enumitem}
\usepackage{makecell}
\usepackage{siunitx}
\usepackage{placeins}
\renewcommand{\baselinestretch}{0.877}
\theoremstyle{definition}
\newtheorem{exmp}{Example}
\DeclareMathOperator{\Expectation}{\mathbb{E}\xspace}
\DeclareMathOperator{\Probability}{\mathbb{P}\xspace}
\usepackage[table]{xcolor}
\usepackage{flushend}

\acrodef{NTN}{Non-Terrestrial Network}
\acrodef{MEC}{Multi-Access Edge Computing}
\acrodef{LEO}{Low Earth Orbit}
\acrodef{DRL}{Deep Reinforcement Learning}
\acrodef{LLM}{Large Language Model}
\acrodef{SLM}{Small Language Model}
\acrodef{SoC}{State of Charge}
\acrodef{ISL}{Inter-Satellite Link}
\acrodef{IoT}{Internet of Things}
\acrodef{KPI}{Key Performance Indicator}

\hyphenation{al-go-rithm allows op-tical net-works semi-conduc-tor ubi-quitous
energy expe-riments satel-lite con-stel-lation orches-tration se-mantic
com-pliance}

\begin{document}
\author{
\IEEEauthorblockN{Octávio B. Nascimento and Carlos A. Astudillo}
\IEEEauthorblockA{%
Institute of Computing, State University of Campinas (UNICAMP),
Campinas, SP, 13083-852, Brazil\\
Emails: o241327@dac.unicamp.br, castudillo@unicamp.br\\
Code: \url{https://github.com/wocn-unicamp/CosmicBeats-Simulator}}
}
\title{\huge Semantic Task Orchestration in LEO-NTN\\
Multi-Access Edge Computing:
A Comparative Study}
\maketitle
    \newcommand{\MYheader}{\smash{\scriptsize
    \hfil\parbox[t][\height][t]{\textwidth}{\centering {}}\hfil\hbox{}}}
    \makeatletter
    \def\ps@IEEEtitlepagestyle{%
    \def\@oddhead{\MYheader}%
    \def\@evenhead{\MYheader}%
    \def\@oddfoot{}%
    \def\@evenfoot{}}
    \makeatother
    \pagestyle{headings}
    \addtolength{\footskip}{0\baselineskip}
    \addtolength{\textheight}{-0.1\baselineskip}
\begin{abstract}
\ac{NTN} \ac{MEC} over \ac{LEO} constellations introduces extreme constraints
on task orchestration: tight energy budgets driven by orbital eclipses,
millisecond-range decision windows, and heterogeneous computing payloads.
Beyond classical resource management, real-world tasks increasingly carry
\emph{semantic constraints}---data sovereignty laws, privacy regulations,
hardware-failure protocols---that pure numerical optimizers cannot interpret.
This work presents a controlled comparative simulation of four orchestration
paradigms over the same 1{,}200-second scenario and identical 69-task workload:
(i)~a greedy heuristic baseline, (ii)~a \ac{DRL} agent (SB3 DQN,
50{,}000~training episodes), (iii)~a \ac{SLM} simulating on-device NPU
inference (Gemma~4 26B, \SI{50}{\milli\second} benchmark latency), and
(iv)~a \ac{LLM} cloud orchestrator (Gemini~3.1 Flash~Lite). Results show that
the LLM achieves 100\% semantic and anomaly compliance---versus 25\% for the
heuristic baseline---at a cost of 5{,}000$\times$ higher energy per decision,
empirically quantifying the compliance-vs-efficiency trade-off for NTN-MEC
system designers.
\end{abstract}
\begin{IEEEkeywords}
NTN-MEC, LEO satellite edge computing, deep reinforcement learning,
large language models, semantic compliance, task orchestration.
\end{IEEEkeywords}
\acresetall
\IEEEpeerreviewmaketitle

\section{Introduction}
\label{sec:intro}

The proliferation of \ac{LEO} constellations has positioned satellite nodes
as candidates for \ac{MEC} infrastructure. When a constellation acts as a
compute fabric, each satellite must independently decide, in real time, whether
to accept an incoming processing task, redirect it to a peer via \ac{ISL}, or
discard it---all subject to strict battery and memory constraints imposed by the
orbital eclipse cycle~\cite{guo2024,gpp2022}.

Beyond classical resource constraints, an emerging class of tasks carries
semantic constraints derived from real-world regulations: data privacy laws
(e.g., GDPR), national data sovereignty requirements, or hardware-level failure
protocols. A compliant orchestrator must parse these constraints from natural
language descriptions and map them onto routing decisions---a capability that
pure numerical optimizers fundamentally lack.

\ac{DRL}, and in particular the value-based Deep Q-Network (DQN) family
introduced by Mnih \emph{et al.}~\cite{mnih2015}, has been applied to satellite
MEC scheduling with promising results. Jiang \emph{et al.}~\cite{jiang2025}
propose a multi-agent DRL framework for on-orbit task scheduling, reporting
improvements of up to 70\% in completion time and 68\% in energy consumption
over unoptimized baselines. Closer to our setting, Casas \emph{et al.}~\cite{casas2026}
address DDQN-based routing in \ac{LEO} constellations and show that agent
performance hinges critically on the reward function, whose design remains a
manual, expertise-intensive process---so much so that they automate it with an
LLM-in-the-loop framework, underscoring the practical cost of reward
engineering. While computationally efficient, such DRL agents operate as
black-boxes over numerical feature vectors, making them unable to interpret
task-level semantic rules without explicit feature engineering.

\acp{LLM} have recently been explored for resource allocation in wireless
systems. Lee and Park~\cite{lee2024} demonstrate that LLMs can perform resource
allocation in wireless networks and explicitly propose a hybrid architecture---a
lightweight DNN for routine decisions combined with an LLM for complex semantic
cases---to compensate for LLM latency limitations, directly motivating the
DRL-vs-LLM trade-off this paper quantifies empirically. Abgaryan
\emph{et al.}~\cite{abgaryan2024} show that LLMs match specialised RL/GNN
approaches in job-shop scheduling benchmarks, confirming that LLM scheduling
capability is not merely a surface-level linguistic effect. The concept of
Edge Intelligence~\cite{deng2020} further motivates embedding compact language
models directly on edge nodes, blurring the boundary between \emph{AI for the
edge} and \emph{AI at the edge}.

For real LEO constellations, Shenoy \emph{et al.}~\cite{shenoy2024} demonstrate
with the CosMAC system that topology-aware scheduling for \ac{IoT} picosat
constellations achieves up to 6.5$\times$ throughput improvement over standard
MAC approaches, confirming that context-sensitive scheduling decisions have
measurable operational impact beyond simulation.

\textbf{Identified gap:} the literature treats DRL, LLM, and Edge Intelligence
as separate threads, each optimising a single numerical metric (time, energy,
throughput). No prior work directly compares all four paradigms---fixed
heuristic, DRL, embedded \ac{SLM}, and centralized LLM---in the same simulated
environment with primary focus on \emph{semantic compliance} rather than
throughput alone. This paper fills that gap with the following contributions:
\begin{enumerate}
  \item A multi-engine comparative simulation framework (4 paradigms,
    deterministic 1{,}200-second window, 69 identical tasks) enabling fair
    paradigm comparison under identical orbital conditions.
  \item A DRL agent trained with SB3 DQN across 50{,}000 synthetic episodes
    using hidden-anomaly reward shaping, achieving 37.5\% anomaly compliance
    without observing semantic rule types---the theoretical maximum for
    binary-flag observations.
  \item Empirical quantification of the compliance-vs-efficiency trade-off:
    LLM achieves 100\% semantic and anomaly compliance at 5{,}000$\times$
    higher energy cost per decision than the heuristic baseline.
  \item An explicit parameter provenance audit (Table~\ref{tab:params})
    distinguishing literature-grounded premises from deliberate design choices.
\end{enumerate}

\section{System Model}
\label{sec:sysmodel}

\begin{figure}[!t]
\centering
\includegraphics[width=\columnwidth]{architecure-diagram.png}
\caption{NTN-MEC orchestration architecture. Three regional \ac{LEO}-\ac{MEC}
nodes (Brazil, Europe, USA) receive a Poisson stream of tasks, each carrying a
region, a RAM requirement, and an optional semantic annotation. A pluggable
orchestration engine---BASELINE, DRL, SLM, or LLM---observes the real-time
satellite state (battery \ac{SoC}, free RAM, solar phase) and maps every task
onto one of four routing actions (SAT-1, SAT-2, SAT-15, or DROP) subject to
battery, memory, and semantic constraints.}
\label{fig:architecture}
\end{figure}

\subsection{Network Architecture}
Fig.~\ref{fig:architecture} depicts the overall orchestration architecture.
A constellation of three \ac{LEO} satellites operates as \ac{MEC} nodes, each
assigned to a fixed geographic region (Brazil, Europe, USA). Satellites carry
heterogeneous computational payloads: 4{,}096~MB RAM each, with battery
capacities of 75~Wh (SAT-1/Brazil), 50~Wh (SAT-2/Europe), and 100~Wh
(SAT-15/USA). The \ac{ISL} is assumed always available (a common assumption for
constellations with optical ISL), so drops reflect exclusively battery/RAM
exhaustion or semantic rule violations---never link failure.

\subsection{Task Arrival and Semantic Anomalies}
Task arrivals follow a Poisson process with rate $\lambda = 4$~tasks/min, the
standard model for aggregated \ac{IoT}/NTN traffic~\cite{guo2024}. Each task
requires 500~MB of RAM held for 60~seconds. A fraction $\alpha = 10\%$ of tasks
carry one of three semantic constraints, assigned uniformly at random:
(i)~\emph{GDPR compliance} (must route exclusively to EUROPE, or discard);
(ii)~\emph{data sovereignty} (must route exclusively to BRAZIL);
(iii)~\emph{critical hardware failure} (must discard regardless of resources).
These three rules constitute the ground-truth reference used to measure
compliance.

\subsection{Battery and Eclipse Model}
Battery dynamics follow a physical orbital eclipse model: solar irradiance
varies sinusoidally with orbital period $T = 90$~min---within the typical LEO
range~\cite{guo2024,gpp2022}---as
$\mathrm{solar}(t) = \sin\!\bigl(2\pi(t+\phi_i)/5{,}400\bigr)$,
where $\phi_i$ is the orbital phase offset of satellite $i$. When
$\mathrm{solar}(t) \le -0.10$ the satellite enters eclipse: battery drains at
a housekeeping rate; otherwise it recharges. This threshold implies
$\approx\!53\%$ eclipse per orbit (46~min), deliberately above the typical
literature range of $\approx$35~min per revolution~\cite{sumanth2019}. The
higher fraction was chosen as a pessimistic stress scenario to generate
observable battery pressure within the 20-minute simulation window. Satellites
refuse new tasks when battery \ac{SoC} falls below 20\% (design choice).

Table~\ref{tab:params} lists all simulation parameters with explicit provenance.

\begin{table}[!t]
\centering
\caption{Simulation Parameters and Provenance}
\label{tab:params}
\footnotesize
\begin{tabular}{p{2.2cm}p{1.5cm}p{3.9cm}}
\toprule
\textbf{Parameter} & \textbf{Value} & \textbf{Source / Rationale} \\
\midrule
Orbital period   & 90 min           & Typical LEO range \cite{guo2024,gpp2022} \\
Eclipse fraction & $\approx$53\%    & Above typical $\approx$35 min \cite{sumanth2019}; pessimistic stress scenario \\
Battery threshold & 20\% SoC        & Design choice \\
Battery capacities & 75/50/100 Wh  & Design choice (heterogeneous HW) \\
SLM latency (NPU) & \SI{50}{\milli\second} & Order-of-magnitude from on-device benchmarks \cite{wang2025} \\
Task arrival     & Poisson, $\lambda$\,=\,4/min & Standard NTN/IoT model \cite{guo2024}; rate is design choice \\
RAM (total / task) & 4{,}096 / 500 MB & Design choice \\
Task hold time   & 60 s             & Design choice \\
Anomaly rate     & 10\%             & Design choice ($\approx$8 anomaly tasks in 69) \\
\bottomrule
\end{tabular}
\end{table}

\section{Orchestration Engines}
\label{sec:engines}

The four engines share a common interface: at each task arrival the engine
receives the task descriptor (region, RAM requirement, semantic annotation if
any) and the real-time satellite states (battery \ac{SoC}, free RAM, solar
phase) and returns one of four actions: route to SAT-1, SAT-2, SAT-15, or DROP.

\subsection{BASELINE (Greedy Heuristic)}
The BASELINE filters candidate satellites by: (1)~region compatibility,
(2)~battery \ac{SoC} $>$ 20\%, and (3)~available RAM $\geq$ required. Among
eligible nodes it selects the highest-battery satellite, breaking ties by
available RAM. The engine executes as a sequence of conditional checks with
negligible latency and has no access to semantic annotations.

\subsection{DRL (SB3 DQN): Problem Formulation}
\label{sec:drl}

We formulate the orchestration decision as a Markov Decision Process (MDP)
$\mathcal{M} = \langle \mathcal{S}, \mathcal{A}, P, R, \gamma \rangle$, following
the problem-formulation methodology of Casas \emph{et al.}~\cite{casas2026} for
RL-based routing in \ac{LEO} constellations, where $\mathcal{S}$ is the state
space, $\mathcal{A}$ the action space, $P$ the (unknown) transition kernel, $R$
the reward function, and $\gamma$ the discount factor. Unlike the multi-hop
packet-forwarding MDP of~\cite{casas2026,jiang2025}, our agent solves a
\emph{single-step admission-and-routing} decision: each task arrival is an
independent episode that terminates immediately after one action ($\gamma = 0$,
a contextual-bandit degenerate MDP). This reflects the operational reality that
the orbital dynamics (battery, eclipse) evolve on a timescale decoupled from the
instantaneous per-task decision, so no intra-episode temporal credit assignment
is required. The agent is \emph{model-free}: it never accesses $P$ nor the
analytic battery/anomaly dynamics of Section~\ref{sec:sysmodel}, learning the
policy purely from sampled transitions.

\subsubsection*{State space $\mathcal{S}$}
Each state $s \in [0,1]^{13}$ concatenates a 3-D one-hot encoding of the task
region, a binary anomaly-present flag $h$, and per-satellite triplets for the
three nodes:
\begin{equation}
s = \bigl[\,\underbrace{r_{\mathrm{US}}, r_{\mathrm{BR}}, r_{\mathrm{EU}}}_{\text{region}},\;
      h,\;
      \{\underbrace{b_i, m_i, c_i}_{\text{sat } i}\}_{i=1}^{3}\,\bigr],
\label{eq:state}
\end{equation}
where $b_i$ is battery \ac{SoC}$/100$, $m_i$ free-RAM fraction $/4096$, and
$c_i \in \{0,1\}$ the solar-charging indicator. Crucially, the state exposes
only the \emph{presence} of an anomaly ($h$), never its \emph{type}. This is a
deliberate scientific invariant: the agent cannot distinguish a GDPR rule from a
data-sovereignty rule from a hardware failure, which structurally upper-bounds
its achievable anomaly compliance (Section~\ref{sec:results}).

\subsubsection*{Action space $\mathcal{A}$}
The action space is discrete with four options,
$\mathcal{A} = \{0,1,2,3\}$, mapping to \{route to SAT-1, route to SAT-2, route
to SAT-15, DROP\}. At inference, the selected action is re-validated against the
live fleet (region match, \ac{SoC}~$>$~20\%, sufficient RAM); an infeasible
choice degrades to DROP, guaranteeing the policy never emits a physically
impossible allocation.

\subsubsection*{Reward function $R(s,a)$}
The reward encodes both resource feasibility and semantic correctness:
\begin{equation}
R(s,a) =
\begin{cases}
+1 + 0.1\,b_a & \text{valid route to eligible sat } a,\\
+1 & \text{DROP on a hardware-failure task},\\
-0.5 & \text{forced DROP (no valid candidate)},\\
-1 & \text{voluntary DROP while a candidate exists},\\
-2 & \text{invalid action (region/\ac{SoC}/RAM fail)},
\end{cases}
\label{eq:reward}
\end{equation}
where $b_a$ is the chosen satellite's \ac{SoC} fraction, so the $+0.1\,b_a$ term
breaks ties toward higher-charged nodes. Because reward design in dynamic
\ac{LEO} settings is itself a hard, expertise-intensive problem---to the point
that recent work automates it with LLMs~\cite{casas2026}---we apply
\emph{reward shaping}: during training only, the hidden anomaly subtype selects
the compliant target region (GDPR$\rightarrow$EUROPE, sovereignty$\rightarrow$BRAZIL)
or mandates a drop (hardware). This is shaping, not information leakage: the
subtype conditions the \emph{reward} but is never part of the observation in
Eq.~\eqref{eq:state}, so at inference the agent must commit to the single best
``blind'' fallback for $h=1$.

\subsubsection*{Dynamics and assumptions}
Training states are drawn i.i.d.: task region uniform over the three regions,
anomaly present with probability $\alpha = 0.10$ (subtype uniform over the three
rules), per-satellite battery $\sim\mathcal{U}(15,100)\%$, free
RAM~$\sim\mathcal{U}(0,4096)$~MB, and solar charging Bernoulli$(0.5)$; task RAM
is fixed at 500~MB. This distribution is designed to cover the operating region
visited by the real SimPy simulation while keeping episodes independent.

\subsubsection*{Algorithm and training}
We use a Deep Q-Network (DQN)~\cite{mnih2015}---a value-based, off-policy,
model-free method---implemented in Stable Baselines3 (SB3). DQN approximates the
action-value $Q(s,a;\theta)$ with a two-layer MLP ($64\times64$) and minimizes
the temporal-difference error toward the bootstrapped target
\begin{equation}
y = r + \gamma \max_{a'} Q(s',a';\theta^{-}),
\label{eq:bellman}
\end{equation}
with target network parameters $\theta^{-}$~\cite{sutton2018,mnih2015}. Since
$\gamma = 0$ in our single-step MDP, Eq.~\eqref{eq:bellman} reduces to
$y = r$, and learning amounts to regressing the immediate reward of
Eq.~\eqref{eq:reward}. The agent is trained for 50{,}000 timesteps with an
$\varepsilon$-greedy schedule and periodic deterministic evaluation
(Table~\ref{tab:dqn}); the best-performing checkpoint is retained. Because SB3
cannot be wired directly into a SimPy discrete-event loop, training and
inference are separated, with the trained model performing deterministic
inference at $\approx$0.475~ms per call. A practical pitfall: \texttt{DQN.load()}
in SB3 silently resets the global Python \texttt{random} seed, corrupting the
deterministic task sequence shared across engines; this was resolved by
restoring the canonical seed immediately after loading.

\begin{table}[!t]
\centering
\caption{DQN Agent Hyperparameters (\texttt{NTNMECEnv})}
\label{tab:dqn}
\footnotesize
\begin{tabular}{p{4.3cm}p{3.3cm}}
\toprule
\textbf{Hyperparameter} & \textbf{Value} \\
\midrule
Policy network        & MLP, $[64, 64]$ \\
Learning rate         & $10^{-3}$ \\
Replay buffer size    & 20{,}000 \\
Batch size            & 128 \\
Discount factor $\gamma$ & 0 (single-step MDP) \\
Target update $\tau$  & 1.0 (hard) \\
Target update interval & 500 steps \\
Train frequency       & every 4 steps \\
Exploration fraction  & 0.30 \\
Final $\varepsilon$   & 0.05 \\
Training timesteps    & 50{,}000 \\
Eval.\ protocol       & 500 deterministic ep.\ / 5{,}000 steps \\
Random seed           & 42 \\
\bottomrule
\end{tabular}
\end{table}

\subsection{Semantic Modeling via Prompting}
\label{sec:prompting}
The SLM and LLM engines replace the numerical feature vector of the DRL agent
with a natural-language prompt, casting orchestration as an instruction-following
task. Fig.~\ref{fig:prompt} shows the full prompt: a role/context header, the
task descriptor (region, RAM, and the raw semantic annotation), the live fleet
state, and an ordered list of four routing rules that map each anomaly onto a
compliant action. This structure mirrors the staged prompt design of
Casas \emph{et al.}~\cite{casas2026}. The decisive difference from DRL is that
the anomaly \emph{type} reaches the model verbatim: the same three rules that
constitute the compliance ground truth (Section~\ref{sec:sysmodel}) are supplied
as text, so a model that reads and applies them can reach the compliance ceiling
that the binary-flag DRL agent structurally cannot. Both models are queried at
\texttt{temperature}=0 for determinism and constrained to emit a single JSON
object; the engine parses it and re-validates the target against the live fleet
before committing.

\begin{figure}[!t]
\centering
\footnotesize
\begin{minipage}{0.97\columnwidth}
\begin{verbatim}
CONTEXT: You are an AI scheduler embedded in a
Low-Earth-Orbit satellite (edge node). ...
Output ONLY valid JSON - no markdown.

TASK:
  region: {region}
  ram_required: {ram} MB
  anomaly: "{anomaly}"

AVAILABLE SATELLITES:
  SAT 1 | region=BRAZIL | battery_pct=64%
        | solar=True | ram_free=3596 MB
  ... (one line per satellite) ...

ROUTING RULES (apply in order):
 1. GDPR/EU privacy anomaly: route
    exclusively to EUROPE. Drop if unavailable.
 2. Sovereignty anomaly: route to task
    country region. Drop if unavailable.
 3. Critical hardware failure anomaly:
    drop the task entirely.
 4. No anomaly: route to a satellite in the
    task region with battery_pct > 20% and
    sufficient RAM.

Output ONLY valid JSON:
{"action": "process|route|drop",
 "target_region": "USA|BRAZIL|EUROPE|null",
 "reason": "one short sentence"}
\end{verbatim}
\end{minipage}
\caption{Natural-language prompt shared by the SLM and LLM engines (SLM variant
shown). Placeholders in braces are filled per task. The LLM variant returns an
explicit \texttt{satellite\_id} instead of an \texttt{action}/\texttt{target\_region}
pair. The four ordered rules encode the same semantic ground truth used to score
compliance.}
\label{fig:prompt}
\end{figure}

\subsection{SLM (On-Device Edge AI)}
The SLM engine simulates on-device NPU inference using the Gemma~4 26B model
via API (parameter count in the class of edge-deployable models). It is modeled
as an \emph{asynchronous, single-lane} accelerator: an arriving task is enqueued
with a ready-time $t_{\mathrm{ready}} = \max(t, t_{\mathrm{busy}}) + \SI{50}{\milli\second}$,
serializing inferences so the NPU processes one decision at a time, and completed
inferences are harvested on later 5-second simulation ticks. The
\SI{50}{\milli\second} figure is the order of magnitude reported for on-device
language-model inference on dedicated NPU hardware~\cite{wang2025}; the real API
call executes out of band. A client-side sliding-window rate limiter (14
requests/min) keeps the engine within the provider quota.

Robust output parsing was the main engineering challenge. Gemma returns its
chain-of-thought and its final answer as separate response \emph{parts}, each
tagged \texttt{thought:true/false}; naively reading the first part produced
$\approx$40\% spurious failures. The parser therefore keeps only non-thought
parts (falling back to concatenation when the final part omits the flag) and, if
standard JSON decoding fails, reconstructs the decision by extracting the
\texttt{action}/\texttt{target\_region} key--value pairs scattered across the
reasoning text. This reduced infrastructure-caused drops from 14 to~2.

\subsection{LLM (Cloud Orchestrator)}
The LLM engine uses Gemini~3.1 Flash~Lite via API, modeling a centralized
ground-station orchestrator, and runs \emph{synchronously} within the decision
loop. It receives the same prompt and rules as the SLM (Fig.~\ref{fig:prompt})
but exploits the provider's native JSON mode (\texttt{responseMimeType}), so no
chain-of-thought parsing is needed. A pre-filter skips the API call entirely
when no satellite in the task region is above the battery threshold, saving cost
on unavoidable drops, and transient rate-limit responses trigger a bounded
exponential back-off. The engine pays the full ground-station round-trip
(\SI{6112}{\milli\second} average, empirically measured), orders of magnitude
above the on-device NPU benchmark---the concrete price of the latency-vs-context
trade-off separating the cloud LLM from the embedded SLM.

\section{Evaluation Metrics}
\label{sec:metrics}

The primary \acp{KPI} are compliance-based, not throughput-based. In the
semantic-constraint setting, a correctly \emph{discarded} task is a success; a
successfully \emph{accepted} but mis-routed task is a failure. Formal
definitions follow.

\begin{itemize}
  \item \textbf{Semantic Compliance Rate (SCR):}
    \begin{equation}
      \mathrm{SCR} = \frac{\text{decisions matching ground-truth rules}}
                          {\text{total tasks}}.
      \label{eq:scr}
    \end{equation}
    For tasks without anomaly, any valid routing to the correct geographic
    region is compliant. For tasks with anomaly, only the rule-prescribed
    action (including mandatory discard) is compliant.

  \item \textbf{Anomaly Compliance Rate (ACR):} identical to SCR, restricted
    to the $\lfloor \alpha N \rfloor$ tasks carrying semantic constraints,
    isolating semantic reasoning capability from basic geographic routing.

  \item \textbf{Task execution rate / drop rate:} fraction of tasks accepted
    (not discarded). A secondary throughput metric; it penalises correct
    discards equally with incorrect ones, making it a poor proxy for
    orchestration quality in the semantic setting.

  \item \textbf{Average decision latency:} mean wall-clock time per
    orchestration call, from conditional check ($\approx$0 for BASELINE) to
    network round-trip (LLM).

  \item \textbf{Energy per decision:} joules consumed by the orchestration
    compute per call, estimated from calibration constants for each paradigm
    (CPU if-else, MLP inference, NPU benchmark, cloud API round-trip). Not a
    direct hardware measurement.

  \item \textbf{Time-weighted RAM utilisation:} integral of
    (RAM in use / total RAM) over simulation time, per satellite, averaged
    across the fleet. By Little's Law ($L = \lambda W$), with
    $\lambda = 4/60$~tasks/s and $W = 60$~s hold time:
    \begin{equation}
      \bar{u} \approx \frac{4 \times 500~\text{MB}}{3 \times 4{,}096~\text{MB}}
               \approx 16\%.
      \label{eq:ram}
    \end{equation}
    Reported as a methodological note in Section~\ref{sec:results}, not as a
    primary KPI, since its expected value is analytically predictable and
    differences among engines are driven by drop rates rather than semantic
    reasoning.
\end{itemize}

\section{Experimental Results}
\label{sec:results}

All four engines process the identical sequence of 69 tasks
under the same 1{,}200-second window and random seed (controlled for fairness).
Table~\ref{tab:kpi} summarises the complete KPI comparison.

\begin{table}[!t]
\centering
\caption{KPI Comparison across Orchestration Paradigms (69 tasks, 1{,}200\,s)}
\label{tab:kpi}
\resizebox{\columnwidth}{!}{%
\begin{tabular}{lcccc}
\toprule
\textbf{KPI} & \textbf{BASELINE} & \textbf{DRL} & \textbf{SLM} & \textbf{LLM} \\
\midrule
SCR (semantic compliance) &
  91.3\% & 92.8\% &
  97.1\% & 100.0\% \\
ACR (anomaly compliance) &
  25.0\% & 37.5\% &
  75.0\% & 100.0\% \\
Task execution rate &
  100.0\% & 88.4\% &
  89.9\% & 95.7\% \\
Drop rate &
  0.0\% & 11.6\% &
  10.1\% & 4.3\% \\
Avg.\ decision latency &
  $\approx$0\,ms & \SI{0.475}{\milli\second} &
  \SI{50}{\milli\second} & \SI{6112}{\milli\second} \\
Total energy &
  \SI{0.069}{\joule} & \SI{0.345}{\joule} &
  \SI{6.9}{\joule} & \SI{345}{\joule} \\
Energy per decision &
  \SI{0.001}{\joule} & \SI{0.005}{\joule} &
  \SI{0.100}{\joule} & \SI{5.0}{\joule} \\
Anomaly tasks (of 69) &
  8 & 8 &
  8 & 8 \\
\bottomrule
\end{tabular}%
}
\end{table}

\subsection{Semantic Compliance}

BASELINE achieves 91.3\% overall SCR but only
25.0\% (2/8) ACR---a result of geographic
coincidence, not semantic reasoning: the engine never observes the presence of
a rule. DRL improves to 92.8\% SCR and
37.5\% ACR (3/8), surpassing BASELINE on both
metrics despite receiving only a binary \emph{anomaly-present} flag. The 37.5\%
ACR ceiling is theoretically justified: of the 8 anomaly tasks, only 3
(mandatory hardware-failure discards) can be correctly identified from a single
binary indicator; the remaining 5 require knowing \emph{which} specific rule
applies (GDPR vs.\ data sovereignty), information the DRL agent structurally
lacks.

SLM reaches 97.1\% SCR and
75.0\% ACR (6/8). The 2 non-compliant cases arose
from API network instability (residual after the thought-parts parsing fix),
not reasoning errors---when the model responded, its decision was correct. LLM
achieves 100.0\% on both metrics, confirming that
reading and applying semantic rules in natural language surpasses both the fixed
heuristic and the purely numerical RL agent, in exact proportion to the semantic
complexity of the task.

The per-anomaly breakdown of the 8 constrained tasks (3~hardware-failure,
3~GDPR, 2~data-sovereignty) makes the compliance ceilings concrete. The
mandatory-discard hardware cases are the \emph{only} ones identifiable from a
binary flag: DRL captures exactly these 3 (its 3 correct drops), which is why it
plateaus at 3/8, while BASELINE---which never discards when a candidate
exists---resolves only the 2 anomaly tasks whose origin region happens to
coincide with the rule-mandated region. The 5 GDPR and sovereignty tasks demand
routing to a \emph{specific} region determined by the rule text; SLM parses this
correctly for 6/8 (missing 2 to transient API errors, not logic), and LLM for
all 8/8. The monotone progression 2$\rightarrow$3$\rightarrow$6$\rightarrow$8
thus tracks precisely how much semantic information each paradigm can access,
from none (geographic coincidence) to full natural-language rule comprehension.

\subsection{Energy Efficiency}

The most dramatic differential is energy per decision (Table~\ref{tab:kpi}).
The LLM-to-BASELINE ratio is $5{,}000$$\times$
(5.0~J vs.\ 0.001~J), reflecting
the full cost of routing each decision through the cloud versus local if-else
logic.

Table~\ref{tab:battery} shows satellite battery state at simulation end under
the BASELINE engine. SAT-2 (EUROPE, smallest 50~Wh capacity) reached
21.8\% \ac{SoC}---only 1.8~percentage points above the
20\% safety threshold---confirming that the pessimistic eclipse model
(Section~\ref{sec:sysmodel}) generates real, measurable energy pressure within
the 20-minute window.

\begin{table}[!t]
\centering
\caption{Final Battery State at $t = 1{,}200$\,s (BASELINE engine)}
\label{tab:battery}
\footnotesize
\begin{tabular}{llcc}
\toprule
\textbf{Satellite} & \textbf{Region} & \textbf{Capacity} & \textbf{Final SoC} \\
\midrule
SAT-1  & Brazil & 75\,Wh  & 61.9\% (Solar) \\
SAT-2  & Europe & 50\,Wh  & 21.8\% $\dagger$ (Eclipse) \\
SAT-15 & USA    & 100\,Wh & 53.4\% (Eclipse) \\
\bottomrule
\multicolumn{4}{l}{\footnotesize$\dagger$ Only 1.8\,pp above 20\% safety threshold.}
\end{tabular}
\end{table}

By Little's Law (Eq.~\ref{eq:ram}), expected steady-state RAM utilisation is
$\approx$16\%. The time-weighted observed values are BASELINE
13.7\%, DRL 12.1\%, SLM
12.3\%, LLM 13.1\%---close
to the theoretical prediction. Differences across engines are driven by drop
rates (fewer accepted tasks $\Rightarrow$ lower average occupancy), confirming
RAM was not the bottleneck: drops reflect semantic and battery constraints, not
queue exhaustion.

\subsection{Latency and Throughput}

\balance
BASELINE and DRL operate in the sub-millisecond range, suitable for any task
arrival rate. SLM operates at the NPU benchmark (\SI{50}{\milli\second}), and
LLM incurs a full round-trip to the ground station
(\SI{6112}{\milli\second} average). From a throughput
perspective, BASELINE accepts all 69 tasks (zero drops); DRL and
SLM each drop $\approx$10\%, and LLM drops only 3 tasks---all
three of which are semantically correct mandatory discards (hardware-failure
protocol), yielding a 100\% effective success rate among tasks the LLM chose
to accept.

\section{Conclusion}
\label{sec:conclusion}

This paper presented the first controlled comparison of four orchestration
paradigms---greedy heuristic, DRL, embedded \ac{SLM}, and centralized
\ac{LLM}---in a semantics-aware NTN-MEC simulation. The central finding is
that LLM achieves 100\% semantic and anomaly
compliance by parsing natural-language rules, while BASELINE and DRL are
limited to 25\% and
38\% anomaly compliance, respectively, by the
structural absence of semantic parsing capability.

Three structural trade-offs define the applicability of each paradigm:
(1)~\emph{Cognitive capacity vs.\ operational cost}: SLM and LLM correctly
handle semantic constraints that BASELINE and DRL cannot identify---at higher
latency and up to $5{,}000$$\times$ higher energy cost per decision.
(2)~\emph{Latency vs.\ context richness}: BASELINE and DRL decide in
microseconds from a minimal state vector; SLM and LLM pay latency proportional
to the reasoning complexity required.
(3)~\emph{Reward-shaped learning vs.\ explicit rules}: reward shaping raised
DRL above BASELINE on both metrics, but the ceiling is structural---without
observing the anomaly type, DRL cannot reach 100\% compliance as zero-shot
LLM/SLM do.

Natural directions for future work include: (a)~a hybrid macro-micro architecture
(LLM for periodic policy reconfiguration; SLM/DRL for high-frequency task-level
decisions), along the lines suggested by Lee and Park~\cite{lee2024};
(b)~longer simulation windows ($\geq$2~hours) to cover multiple complete eclipse
cycles and allow DRL convergence on the full state distribution; (c)~higher-load
stress scenarios ($\lambda \geq 12$~tasks/min) to saturate RAM and decouple
queuing effects from semantic compliance effects.

\FloatBarrier
\begin{thebibliography}{10}

\bibitem{guo2024}
Z.~Guo, Q.~Chen, and W.~Meng,
``Service-Oriented AoI Modeling and Analysis for Non-Terrestrial Networks,''
arXiv:2408.17051, 2024.

\bibitem{gpp2022}
3GPP Release~17,
``Study on New Radio (NR) to Support Non-Terrestrial Networks,''
TR~38.821, 2022.

\bibitem{jiang2025}
Q.~Jiang, P.~Han, X.~Xin, and K.~Chen,
``Deep Reinforcement Learning and Edge Computing for Multisatellite On-Orbit
Task Scheduling,''
\emph{IEEE Trans.\ Aerosp.\ Electron.\ Syst.}, 2025.
DOI:~10.1109/TAES.2025.3583276.

\bibitem{lee2024}
W.~Lee and J.~Park,
``LLM-Empowered Resource Allocation in Wireless Communications Systems,''
arXiv:2408.02944, Aug.\ 2024.

\bibitem{abgaryan2024}
H.~Abgaryan, A.~Harutyunyan, and T.~Cazenave,
``LLMs can Schedule,''
arXiv:2408.06993, Aug.\ 2024.

\bibitem{deng2020}
S.~Deng \emph{et al.},
``Edge Intelligence: The Confluence of Edge Computing and Artificial
Intelligence,''
\emph{IEEE Internet Things J.}, 2020.
arXiv:1909.00560.

\bibitem{shenoy2024}
J.~Shenoy \emph{et al.},
``CosMAC: Constellation-Aware Medium Access and Scheduling for IoT
Satellites,''
\emph{Proc.\ ACM MobiCom}, 2024.
DOI:~10.1145/3636534.3690657.

\bibitem{sumanth2019}
R.~M.~Sumanth,
``Computation of Eclipse Time for Low-Earth Orbiting Small Satellites,''
\emph{Int.\ J.\ Aviation, Aeronautics, and Aerospace}, vol.~6, no.~5, 2019.
DOI:~10.15394/ijaaa.2019.1412.

\bibitem{wang2025}
H.~Wang \emph{et al.},
``lm-Meter: Unveiling Runtime Inference Latency for On-Device Language
Models,''
arXiv:2510.06126, 2025.

\bibitem{sutton2018}
R.~S.~Sutton and A.~G.~Barto,
\emph{Reinforcement Learning: An Introduction}, 2nd~ed.
MIT Press, 2018.

\bibitem{mnih2015}
V.~Mnih \emph{et al.},
``Human-Level Control through Deep Reinforcement Learning,''
\emph{Nature}, vol.~518, no.~7540, pp.~529--533, 2015.
DOI:~10.1038/nature14236.

\bibitem{casas2026}
W.~P.~Casas, N.~L.~S.~da~Fonseca, and C.~A.~Astudillo,
``LLM-Driven Automated Reward Design for Reinforcement Learning-Based Routing
in LEO Satellite Networks,''
in \emph{Proc.\ IEEE Global Communications Conf.\ (GLOBECOM)}, 2026.

\end{thebibliography}

\end{document}